\section{Good polynomial labels on the original automorphic rows}

% BEGIN G85M1
\paragraph{G85M1. A simultaneous finite type on the selected row.}
Use the original layer-point guards and fix a member $D=w.D(Z)$,
its nonempty member data, coefficient pin and chosen row reading
$Y=t.r.Y$. Write $r$ for the row rank. For every $\sigma$ the
representation $\pi=Y(\sigma)$ occurs by the defining test on the
row. Thus for one of the two levels there is a nonzero multiplicity
vector $T$ in the intersection of the good subspace with the
isotypic subspace for $\bar\pi$. The conjugation here comes from
the row occurrence test. We describe the effect on all the finite
labels using this same vector.

At a designated split place above $g_k$, identify its completion with
$\mathbb Q_{g_k}$. The two unitary coordinates of a local lift are
the chosen matrix and its inverse transpose in paired frames. A lift
of any matrix in $\mathrm{GL}_r(\mathbb F_{g_k})$ exists: lift its
entries to the integers, where the determinant is a unit, and use
identity at all other places. Changes from the principal frame to
these frames are conjugations; the finitely many changes are the
components of one finite adele $c$. The level at either of the two
rows contains the kernel of reduction at each type place, acting
with trivial character there. After the transport by $c$, the
principal congruence groups thus fix $T$. On their fixed vectors the
lifts give an action of the product $G$ of the finite groups: two
lifts of one matrix differ by the kernel and different places
commute. The row pins equate the coefficients of the finite
operator $e_0$ in these coordinates with those of the geometric
good projector. Applying the equality term by term to $T$, whose
values lie in the prescribed level type, gives $e_0T=T$.

The isotypic subspace is a sum of copies of the irreducible
$\bar\pi$, hence a direct sum of such copies after choosing a
maximal independent family. Some projection onto one copy detects
$T$. Compose with right translation by $c^{-1}$ to intertwine the
conjugated action with ordinary right translation. Its image $x$
is nonzero, fixed by the principal congruence groups and by $e_0$.
By P84C1 there is a good tuple $f$ for which
$y=(\prod_{kv}e_{f_{kv}})x\ne0$. Each local factor fixes $y$.
Conjugate this vector and pass to the local representation
$\pi_{kv}$. More explicitly, the translates away from that place
of an embedded copy of $\pi_{kv}$ span the global irreducible:
their sum is stable under the local group and the complementary
group and is nonzero. They are simple copies for the local action.
Projection onto one of their direct summands detects $\bar y$.
It intertwines the compact group, hence gives a nonzero
principal-congruence fixed vector in $\pi_{kv}$ on which
$\overline{e_{f_{kv}}}$ acts as identity. The local projection may
vary with $k,v$; the tuple $f$ was chosen only once.
% END G85M1

% BEGIN G85M2
\paragraph{G85M2. Matching a finite series with the specified tame generator.}
At one of these places put $g=g_k$, $K=\mathrm{GL}_d(\mathbb Z_g)$
in rank $d$, and $K_1=\ker(K\to\mathrm{GL}_d(\mathbb F_g))$.
Use the original weak local correspondence and the row's
Weil--Deligne comparison. We show that any finite series $P$ acting
nontrivially on $\pi_{kv}^{K_1}$ has inertia polynomial
$\prod_{x\in\mathrm{roots}(P)}(X-\mathrm{Teich}(x))$ at
the named generator $s$. Roots always count multiplicities.

First take a supercuspidal support factor with nonzero
$K_1$-fixed vectors. The depth-zero classification used in A72F7
realizes it as compact induction from the center times $K$ of an
extension of an inflated finite cuspidal type. Only the trivial
Cartan double coset modulo the center contributes $K_1$-invariants:
for unequal elementary divisors choose a strict gap; the
intersection of the conjugated $K_1$ with $K$ contains modulo $g$
the whole rectangular unipotent radical across that gap, since the
entry scalings there remove at least one factor of $g$.
A finite cuspidal type has zero invariants under such a radical.
Thus its finite series is its Green cuspidal polynomial, a monic
irreducible $A$ of degree $d$. This argument includes unramified
twists; in degree one the same assertion reads the unit character.

Choose a root $a$ of $A$. In $\mathbb F_{g^d}^\times$ use the
member's designated generator $h_d$. The character
$\theta(h_d^j)=\mathrm{Teich}(a)^j$ is well defined and regular,
since the $g$-orbit of $a$ has size $d$. Green's elliptic character
formula is the signed sum of the $d$ Frobenius conjugates of this
character, exactly the regularity test in the type law of the weak
correspondence. To evaluate that law use $E'=E(\mu_{g^d-1})$:
at the chosen split place its completion is unramified of degree $d$
because $g$ has order $d$ modulo $g^d-1$ (in degree one the
completion is unchanged). Thus $g$ is a uniformizer there and
$s$ acts trivially on $E'$. The member's generator pin says that
for $z^{g^d-1}=g$ the residue of $s(z)/z$ is $h_d$, under the
specified residue identification. This is the residue of the local
Artin unit of $s$ with the convention of the type law. To check
the direction, put $Q=g^d$. For a residue unit $u$ choose
Frobenius lifts fixing $z$ and $z u^{1/(Q-1)}$, respectively,
by taking inverse Artin uniformizers in the corresponding totally
ramified Eisenstein extensions and descending by the norm.
Their quotient (inverse of the second times the first) represents
the Artin class of $u$: the two uniformizers $g,gu$ have norm
representatives $(-1)^Qg,(-1)^Qgu$ and the inverse Frobenius
normalization takes their ratio. On $z$ the quotient has multiplier
of residue $u$, since arithmetic Frobenius raises the residue
of $u^{1/(Q-1)}$ to its $Q$th power. Principal units are
$(Q-1)$st powers by Hensel, hence do not change this tame
quotient. Since $\mu_{Q-1}\subset E'$, the multiplier is a
character and depends only on the abelianized inertia element.
The local Artin unit of $s$ therefore has the asserted residue.
The type law evaluated on this unit now gives the roots
$\theta(h_d^{g^j})=\mathrm{Teich}(a^{g^j})$. For the
rank-restricted native law its numerical guard holds throughout:
$d\le r\le n_V$, the host has $q,d_{\rm host}\ge n_V$ and
$16\mid Q_S$, and $Q_S q d_{\rm host}\mid g-1$.

For general $\pi_{kv}$ take a Zelevinsky supercuspidal support
list whose parabolic induction has this representation as a
subquotient. Compact averaging makes $K_1$-invariants exact.
Iwasawa decomposition restricts the induction to induction from
the integral parabolic. Taking $K_1$-invariants identifies it with
finite parabolic induction of the $K_{1,j}$-invariants of the
factors: functions descend to $K/K_1$, and evaluating at coset
representatives imposes exactly the Levi congruence invariance.
The modulus character is trivial on the integral parabolic.
The cuspidal labels occurring above therefore multiply to $P$
by finite Harish-Chandra support (Green's series); all ranks are
at most $r$. On the parameter side the segment law takes direct
sums of the support factors with the indicated unramified twists;
these twists are trivial on inertia. Multiplying the preceding
cuspidal formulas proves the claimed polynomial.
% END G85M2

% BEGIN G85M3
\paragraph{G85M3. All original rows pass both goodness tests.}
For a monic $P$ with nonzero constant term put
$P^*(X)=X^{\deg P}P(X^{-1})/P(0)$. Conjugating a finite
cuspidal character sends each Green parameter $a$ to $a^{-1}$:
its elliptic character sum has real sign and its values in each
summand are roots of unity. Finite parabolic induction commutes
with complex conjugation. Thus conjugation carries the entire
finite series of $P$ to that of $P^*$, including every simple
constituent. The central character idempotent formula
$e_\lambda=\dim(\lambda)|G|^{-1}\sum_h
\mathrm{tr}(\lambda(h^{-1}))h$ gives
$\overline{e_P}=e_{P^*}$. G85M1 and G85M2 therefore give
the lifted root polynomial of $f_{kv}^*$ on the actual local
parameter. The chosen member's $\mathrm{row}$ and
$\mathrm{wdT}$ identify this local parameter with
$\mathrm{locR}(\sigma,w_{kv})$ in the original inertia-polynomial
coordinate. The dimension law gives $r$ there; the row label
has degree $r$, so $\mathrm{inertiaCharpoly\_type}$ and
$\mathrm{row.lab\_eq}$ identify the same characteristic
polynomial with the lift of $\mathrm{lab}(k,\sigma,v)$.
Here the finite series and Teichm\"uller map are the member's
actual ones by $t.\mathrm{grp.ser}$ and $t.\mathrm{grp.teich}$.
Teichm\"uller lifting is injective on nonzero roots. The lift also sends zero to zero: otherwise the identity
$\mathrm{Teich}(0)\mathrm{Teich}(x)=\mathrm{Teich}(0)$ would
force every unit to have the same image. Units map to units.
The full lifted root multiset therefore determines the original monic
polynomial, and $\mathrm{lab}(k,\sigma,v)=f_{kv}^*$.

Reciprocation inverts the roots, preserves squarefreeness and
gives a degree-preserving bijection on monic divisors, taking
irreducible divisors to irreducible divisors. Inversion commutes
with projection to the prime-to-$\ell$ part, so preserves the
root separation tests and exchanges the two tests against $\chi$
and $\chi^{-1}$. The places witnessing missing divisor degrees
and a large irreducible prime degree work without change.
Both original goodness predicates for $f$ therefore give both
for its reciprocal tuple. This proves $D.\mathrm{nd.goodRow}(\sigma)$
for every original $\sigma$. The quantification is on the selected
member itself: the row representation is transported by the
chosen row equality, and $D=\mathrm{withLoc}(c.a.\mathrm{nt}.L)$
identifies its labels with the ones just computed.
This supplies $\mathrm{GaloisPreconditions.good}$.
% END G85M3
